\subsection{Requerimientos, supuestos y restricciones}

\subsubsection{Requerimientos}

Este proyecto cuenta con una serie de requerimientos a cumplir, empezando por las fechas importantes de entrega de avances así como la entrega final, todo esto siguiendo la metodología CRISP-DM, programado para las siguientes fechas:

\begin{itemize}
  \item Avance sobre el apartado de ``Business Understanding'' programado para el 7 de octubre.
  \item Avance sobre ``Data Understanding'' para el día 20 de octubre.
  \item Avance en la sección ``Data Preparation'' estipulada para el 4 de noviembre.
  \item Como cuarto avance se tiene a la parte de ``Modelling'' en el día 15 de noviembre.
  \item Finalmente como avance se tiene la sección ``Evaluation'' predicha para el 20 de noviembre.
  \item La entrega final de todo el proyecto sería el 2 de diciembre.
\end{itemize}

El segundo requerimiento a tener en cuenta son los datos, cómo es que vienen pero enfocado en su comprensión y calidad, en este caso no se profundiza mucho en ello ya que esto es parte de la sección de ``Data Understanding'' siguiendo con la metodología CRISP-DM. Por lo que solo nos centramos en explicar lo justo y necesario. Los datos que tenemos son imágenes en formato HEIC y JPG que se tomaron en las dos visitas al CAETEC (29 de septiembre y 14 de octubre de 2026), siguiendo con una vista posterior a la vaca, alineada a ella y desde más arriba; también fueron laterales de lado izquierdo y derecho, esto para seguir con un estándar en todas las vacas fotografiadas. Buscamos tomarlas de la misma manera para reducir la variabilidad que pueda hacer que falle el modelo, aunque al usar teléfonos distintos quedaron con varias resoluciones (véase Comprensión de los datos).

Para el apartado de seguridad y legal, tenemos acceso a los datos y los podemos utilizar siempre y cuando no los introduzcamos en sitios públicos en donde otras personas ajenas a la institución puedan utilizarlos para sus intereses personales, al no estar tratando con datos de personas no existe un riesgo de que alguien salga afectado al usar sus datos, pero sí estamos tratando con el trabajo de recolección de otros, por ello es importante cuidar el uso que le damos ya que no nos pertenecen.

\subsubsection{Supuestos}

\begin{itemize}
  \item Todas las vacas se encuentran en un rango de BCS de 2.0 a 3.75.
  \item El BCS dicho por los expertos será tomado como lo correcto.
  \item Entre más leche producen las vacas mejor.
\end{itemize}

\subsubsection{Restricciones}

\begin{itemize}
  \item Solo se contó con dos visitas al CAETEC (29 de septiembre y 14 de octubre de 2026) para conocer, entender y tomar fotografías.
  \item No todas las imágenes cuentan con la misma iluminación, puede ser peor o mejor.
  \item Las fotos se toman con los teléfonos de los integrantes de los equipos, que no son del mismo modelo (iPhone 17 Pro Max, iPhone 14, iPhone Air, iPhone XR, Samsung S22 y Nothing Phone (2a)).
  \item Se necesita encontrar una computadora con GPU para poder entrenar el modelo.
\end{itemize}
